AIsafePy reduces risk in measurable, well-scoped ways. It does not
provide provable security. This section lists what the library does
not do, the assumptions it makes, and how deployers should think
about residual risk.

\subsection{Limits of capability-based IFC}

CaMeL-style IFC depends on an assumption that the privileged LLM
(P-LLM) emits only the constrained action DSL. If the planner gets
jailbroken into emitting raw natural-language reasoning that the
host code then evaluates, the guarantees degrade to whatever the
host's parser is willing to accept. AIsafePy's adapter layer
protects against this in a layered way: tool arguments are wrapped
in \texttt{Tainted} values at the SDK boundary regardless of how
the planner reasoned, and the policy's
\texttt{control\_flow\_integrity} requirement still rejects calls
that came from UNTRUSTED context. But there are legitimate
adversarial workflows where the LLM produces tool arguments that,
on the surface, appear TRUSTED. Those get caught by deny-if
predicates, by capability mismatches, or, in mediated mode, by the
human approver.

CaMeL reports about $2.7\times$ token overhead. This is a real cost.
Deployers should plan for it and instrument utility before and
after enabling \texttt{aisafepy.flow}.

Even system-level defenses can be bypassed by obfuscation,
encoded payloads, or attacks against the planner's parsing rules.
AIsafePy.flow reduces the attack surface significantly but does
not eliminate it. The honest framing is ``architectural defense
in depth,'' not ``injection is solved.''

\subsection{Limits of streaming forecasters}

Forecasting-based streaming detectors require either Monte-Carlo
rollouts or token-level supervision at training time. The v0.1
release ships pre-trained forecaster wrappers and a recipe for
retraining, not a fully automated training pipeline. Practitioners
who want to retrain on their domain-specific unsafe distribution
have to follow the recipe documented in the \texttt{forecasters}
module docstring.

Activation probes are model-specific. A linear probe trained on
Llama-3-8B does not transfer to Qwen3-8B. We document this in the
package and ship a small set of pre-trained probes per popular
open-weights family in a separate model zoo.

\subsection{Limits of the eval-to-guardrail compiler}

The compiler's output is only as good as the cluster boundaries it
finds. HDBSCAN does well on most realistic red-team logs we have
tested, but adversarially designed attacks that spread across many
small clusters can dilute the signal. The current implementation
has a \texttt{min\_attack\_success\_rate} threshold below which
clusters are dropped; tuning this is a deployment decision.

Steering-vector deployment works for self-hosted models only. For
hosted APIs the compiler falls back to classifier, regex, and
deliberative-case targets. The mean-difference baseline we ship
matches the simplest contrastive activation addition recipe
\citep{panickssery2023cad}; users running conditional or
representation-engineering variants should substitute their own
training loop.

The canary deployment checks the false-positive rate against a
labeled benign baseline. If the baseline is unrepresentative of
production traffic, the rollback signal is unreliable. We recommend
constructing the baseline from logged production traces specifically
marked as benign.

\subsection{Adversarial false positives}

Short adversarial prompts can drive content classifiers to block a
large fraction of benign traffic, producing a denial-of-service
vector that targets the guardrail rather than the underlying model.
AIsafePy.stream mitigates this with three structural defenses: an
input-length cap on every Tier-1 regex guard, regex-catastrophic-
backtracking protection through the same cap, and a planned per-IP
false-positive-rate monitor in the upcoming v0.2 release. The
underlying classifier vulnerability is not eliminated. We recommend
deployers pair the cascade with rate limiting at the request edge.

\subsection{What AIsafePy does not solve}

AIsafePy does not solve sleeper-agent or deceptive-alignment
problems. The right frame is ``defense in depth with measurable
risk reduction.''

Compliance certification (OWASP top-10 for LLMs, NIST AI RMF, EU AI
Act) is downstream of the technical primitives. The library
provides audit-log artifacts that help compliance teams, but
certification itself belongs to a separate process.

The library does not solve hosted-API governance. If a provider
rotates a model behind the same API endpoint, AIsafePy can detect
behavioral drift via the production-traces adapter but cannot
prevent the rotation.

\subsection{Composition with other defenses}

AIsafePy is designed to compose, not to compete. The recommended
deployment pairs \texttt{flow} (architectural) with \texttt{stream}
(detective) and a periodic \texttt{adapt} run (continuous).
External classifiers from \texttt{contrib/} sit inside
\texttt{stream}'s tiers. PyRIT continues to drive \texttt{adapt}.
Inspect AI continues to drive evaluation. The library's value is
the connective tissue, not a replacement for any specific upstream
tool.
